%! AP Statistics — Unit 8, Lesson 8.1 — Exit Ticket (Answer Key)
\documentclass[10pt]{article}
\usepackage{apstats-article}
\usepackage{apstats-key}

\begin{document}

\pageheader{Unit 8, Lesson 8.1}{Exit Ticket --- Answer Key}
\namedateperiod

\vspace{0.1in}
A spinner has 4 equal sections labeled A, B, C, and D. It is spun 80 times with the
following results:

\vspace{0.08in}
\renewcommand{\arraystretch}{1.3}
\begin{tabularx}{0.6\linewidth}{Xcccc}
    \textbf{Outcome} & A & B & C & D \\
    \hline
    \textbf{Observed} & 25 & 18 & 22 & 15 \\
\end{tabularx}

\vspace{0.15in}

\begin{enumerate}

  \item If the spinner is fair, how many times would you expect each outcome to occur?

        \vspace{0.05in}
        Expected count for each outcome $= $ \ans{$80 \div 4 = 20$ each}

        \vspace{0.1in}

  \item Which outcome shows the largest difference from expected? How large is that
        difference?

        \ansline{Outcome D: observed 15, expected 20, difference $= 15 - 20 = -5$.}

  \item Just by looking at these differences, do you think the spinner is unfair?
        What additional information would make you more confident in your answer?

        \ansline{Answers vary --- accept reasonable judgment. The differences are small relative to an expected count of 20 (the largest is only $-5$). A chi-square test would let us compute a $p$-value and quantify how unusual these results are if the spinner is actually fair.}

        \begin{teachernote}
            Students should recognize that the gaps here are not dramatic
            (all within $\pm 5$ of 20). The intent is for students to articulate
            that informal judgment is limited and that a formal test provides the
            $p$-value needed for a decision. Accept any answer that acknowledges
            uncertainty and mentions that a test is needed.
        \end{teachernote}

\end{enumerate}

\end{document}
